\section{Robustness}
\label{sec:robustness}

\subsection{Is there one best policy across playstyles?}
\label{sec:oneBest}

This is the question the project set out to answer, and it deserves a direct
answer rather than a table. The honest one has two halves, and the useful
observation is that they correspond to the two halves of the agent.

\paragraph{The inference half is playstyle-independent, and provably so.}
Constraints (C1)--(C5) are logical consequences of the rules. ``The asker holds
another card of that half-suit'' is true whether the asker is a world champion,
a novice, or a uniform random control; so is ``a miss proves the target lacks
the card'', and so is every capacity identity. Nothing in the exact posterior of
\S\ref{sec:belief} is fitted to an opponent, and no opponent can make it wrong.
The same is true of Theorem~\ref{thm:locked}: a half-suit held by one team cannot
be asked in by the other, whatever the other team's style. So a genuinely
universal core exists, and it is the part of FishBot v0.4 that the ablations
identify as carrying the most value.

\paragraph{The decision half cannot be universally optimal, and we should not
claim it is.} Choosing among asks trades off quantities --- how dangerous it is
to hand this opponent the turn, how much an ask reveals, when to cash a
half-suit --- whose values depend on what the other five players do next. In game
theoretic terms, the best response to a policy is a function of that policy, so
no single deterministic policy is simultaneously a best response to all of them.
The strongest defensible target is not ``best against everyone'' but
``minimally exploitable'', and that is a team-maxmin-with-correlation equilibrium
which we do not compute. Our fitting objective is a deliberate approximation of
it: we maximise the worst-case win rate over a panel rather than the average
(\eqref{eq:softmin}), which is a max-min over a finite, hand-built opponent set
rather than over all policies.

\paragraph{What the evidence supports.} Within the population we can construct,
one policy \emph{is} best against every member simultaneously --- FishBot v0.4's
worst result against any of the eight is \numWorstCase\%, and it is ahead of
every one of them. That is the strongest empirical form of the claim available
without an exploitability bound. Against a purpose-built adversary the picture
changes: an exploiter fitted specifically against it reaches only \numLBRFour\%
(\S\ref{sec:oneBest} is answered in the affirmative for this population, and
\S\ref{sec:lbr} qualifies it). The two results are not in
tension; they say that the agent is uniformly strong against the styles that
exist and remains exploitable by a style built to beat it, which is the expected
state of affairs for any policy short of an equilibrium.

\subsection{Across playstyles}

The fitting objective \eqref{eq:softmin} targets the worst opponent rather than
the average one, and Table~\ref{tab:h2h} reports the whole profile rather than a
row average. Two of the policies in that table --- the misdirection artist and
the random control --- were deliberately excluded from the fitting panel and are
therefore held-out opponents in the strict sense.

\subsection{Across rule dialects}

Because the v0.4 engine implements the rule differences of \S\ref{sec:rules} as
switches, the same frozen policy can be replayed under other dialects.
Table~\ref{tab:rules} reports three: the v0.3 dialect in which declarations are
restricted to the declarer's own turn; the 48-card, eight-half-suit Literature
form; and the full rules with out-of-turn declarations disabled but everything
else intact.

\begin{table}[ht]
\centering
\caption{The frozen FishBot v0.4 configuration replayed under other rule
dialects.}
\label{tab:rules}
\small
\begin{tabular}{llrr}
\toprule
Rule dialect & Opponent & Win rate & 95\% CI \\
\midrule
v0.3 dialect (turn-only declarations) & FishBot v0.3 & 64.46\% & 62.6--66.4 \\
48-card, eight half-suits & FishBot v0.3 & 72.54\% & 70.8--74.2 \\
No out-of-turn declarations & FishBot v0.3 & 67.71\% & 65.9--69.5 \\
v0.3 dialect (turn-only declarations) & Turn-starvation lockout & 77.25\% & 75.6--78.9 \\
48-card, eight half-suits & Turn-starvation lockout & 74.79\% & 73.2--76.4 \\
No out-of-turn declarations & Turn-starvation lockout & 80.17\% & 78.5--81.8 \\
v0.3 dialect (turn-only declarations) & Posterior detective & 76.21\% & 74.4--78.0 \\
48-card, eight half-suits & Posterior detective & 72.21\% & 70.5--73.9 \\
No out-of-turn declarations & Posterior detective & 74.75\% & 73.0--76.5 \\
\bottomrule
\end{tabular}
\end{table}

\subsection{Against an opponent that has studied it}
\label{sec:lbr}

Head-to-head results against a fixed population say nothing about how a
well-informed adversary would fare. Because exact best-response computation is
out of reach at this game size, we estimate a \emph{local} best response: we
freeze FishBot v0.4 and run the same fitting machinery with that frozen policy as
the entire opponent panel, so the optimiser's only objective is to beat it. The
resulting win rate is an estimate of how much an adversary who can study the bot
and search the same policy class can extract. We repeat the procedure against
FishBot v0.3 for comparison, since an exploitability number is only interpretable
relative to something.

The outcome is the most interesting single result in this paper. An adversary
fitted specifically against the frozen FishBot v0.4, inside the same policy class
and with the same optimiser and budget, reaches \numLBRFour\% on a bank neither
the exploiter nor the target had seen --- an interval that straddles even money,
so the probe simply fails to establish any advantage over its target. The identical procedure applied to FishBot v0.3 reaches \numLBRThree\%.

The gap is what matters. A fixed opponent pool cannot distinguish a policy that
is strong from one that is strong \emph{and hard to counter}, and by this measure
the two predecessors are very different objects: v0.3 is comfortably exploitable
by a purpose-built adversary, and v0.4, within the reach of this probe, is not.

The precise reading is worth stating, because it is stronger than ``hard to
beat''. The exploiter searches the same parametric family that FishBot v0.4
belongs to, with the same optimiser and budget, and its objective is solely to
beat v0.4. That it converges to parity means no member of that family, found by
that search, is a profitable deviation --- which is the empirical signature of an
approximate symmetric equilibrium \emph{of the restricted game} in which both
teams are confined to this policy class. It is emphatically not an equilibrium of
Fish: the class is a twenty-feature linear score with a dozen decision knobs, and
the true strategy space is vastly larger. But it does say that the fitting
procedure has run out of room against itself, which is the point at which a
larger policy class, not more fitting, is what the next version needs.

Three caveats keep this from being an exploitability bound. The probe searches
one parametric policy class with one optimiser for a fixed budget, so it is a
lower bound and a differently-shaped adversary could do better. The interval
straddles $50\%$, so ``cannot beat it'' is a failure to establish an advantage
rather than a demonstrated absence of one. And a policy playing a near-copy of
itself is a mirror match, which is precisely where
\S\ref{sec:termination}'s deadlock pressure bites hardest; the residual rate at
which that probe's games reach the action cap is reported in the artifact and is
materially higher than anywhere else in this study.

\begin{table}[ht]
\centering
\caption{Local best-response probe. Each row fits a fresh policy, inside the same
policy class and with the same optimiser, whose only objective is to beat the
frozen policy named. The win rate is what that adversary achieved; lower is
better for the frozen policy.}
\label{tab:lbr}
\small
\begin{tabular}{lrrr}
\toprule
Frozen target & Response win rate & 95\% CI & Mean half-suits conceded \\
\midrule
\vfast{} & 51.19\% & 49.67--52.72 & 4.550 \\
\vthree{} & 77.31\% & 75.97--78.67 & 5.547 \\
Posterior detective & 76.47\% & 75.11--77.81 & 5.471 \\
\bottomrule
\end{tabular}
\end{table}

Table~\ref{tab:lbr} gives the three probes. The probe is a lower bound on
exploitability, not an upper one: a better
optimiser, a longer budget, or a policy class outside the one searched could all
do better. It is reported because a head-to-head result against a fixed
population says nothing at all about an adversary who has studied the agent, and
because an exploitability figure is only interpretable next to another one.

